% Options for packages loaded elsewhere
\PassOptionsToPackage{unicode}{hyperref}
\PassOptionsToPackage{hyphens}{url}
%
\documentclass[
  openany]{book}
\usepackage{amsmath,amssymb}
\usepackage{iftex}
\ifPDFTeX
  \usepackage[T1]{fontenc}
  \usepackage[utf8]{inputenc}
  \usepackage{textcomp} % provide euro and other symbols
\else % if luatex or xetex
  \usepackage{unicode-math} % this also loads fontspec
  \defaultfontfeatures{Scale=MatchLowercase}
  \defaultfontfeatures[\rmfamily]{Ligatures=TeX,Scale=1}
\fi
\usepackage{lmodern}
\ifPDFTeX\else
  % xetex/luatex font selection
\fi
% Use upquote if available, for straight quotes in verbatim environments
\IfFileExists{upquote.sty}{\usepackage{upquote}}{}
\IfFileExists{microtype.sty}{% use microtype if available
  \usepackage[]{microtype}
  \UseMicrotypeSet[protrusion]{basicmath} % disable protrusion for tt fonts
}{}
\makeatletter
\@ifundefined{KOMAClassName}{% if non-KOMA class
  \IfFileExists{parskip.sty}{%
    \usepackage{parskip}
  }{% else
    \setlength{\parindent}{0pt}
    \setlength{\parskip}{6pt plus 2pt minus 1pt}}
}{% if KOMA class
  \KOMAoptions{parskip=half}}
\makeatother
\usepackage{xcolor}
\usepackage{longtable,booktabs,array}
\usepackage{calc} % for calculating minipage widths
% Correct order of tables after \paragraph or \subparagraph
\usepackage{etoolbox}
\makeatletter
\patchcmd\longtable{\par}{\if@noskipsec\mbox{}\fi\par}{}{}
\makeatother
% Allow footnotes in longtable head/foot
\IfFileExists{footnotehyper.sty}{\usepackage{footnotehyper}}{\usepackage{footnote}}
\makesavenoteenv{longtable}
\usepackage{graphicx}
\makeatletter
\def\maxwidth{\ifdim\Gin@nat@width>\linewidth\linewidth\else\Gin@nat@width\fi}
\def\maxheight{\ifdim\Gin@nat@height>\textheight\textheight\else\Gin@nat@height\fi}
\makeatother
% Scale images if necessary, so that they will not overflow the page
% margins by default, and it is still possible to overwrite the defaults
% using explicit options in \includegraphics[width, height, ...]{}
\setkeys{Gin}{width=\maxwidth,height=\maxheight,keepaspectratio}
% Set default figure placement to htbp
\makeatletter
\def\fps@figure{htbp}
\makeatother
\setlength{\emergencystretch}{3em} % prevent overfull lines
\providecommand{\tightlist}{%
  \setlength{\itemsep}{0pt}\setlength{\parskip}{0pt}}
\setcounter{secnumdepth}{5}
\ifLuaTeX
  \usepackage{selnolig}  % disable illegal ligatures
\fi
\usepackage{bookmark}
\IfFileExists{xurl.sty}{\usepackage{xurl}}{} % add URL line breaks if available
\urlstyle{same}
\hypersetup{
  pdftitle={Introduction to RDM},
  pdfauthor={Louise Gillis},
  hidelinks,
  pdfcreator={LaTeX via pandoc}}

\title{Introduction to RDM}
\author{Louise Gillis}
\date{}

\begin{document}
\maketitle

{
\setcounter{tocdepth}{1}
\tableofcontents
}
\chapter{Basic RDM}\label{basic-rdm}

Some intro material on RDM

\chapter{Introduction}\label{introduction}

What is \textbf{research data}? Research data are measurements, recordings, records, or observations about the world collected by researchers with minimal contextual interpretation. They can show up in just about any format. Numbers, notes, images, films, designs, algorithms, and diagrams can all be research data. The takeaway? Research data isn't the exclusive purview of bench scientists - anyone who's engaged in the research process is a collector of, and manager of, research data.

\textbf{Research Data Management} describe the collective process of structuring, organizing, maintaining, and caring for research data. It's straight forward (but not always easy!) steps like:

\begin{itemize}
\tightlist
\item
  creating a readme file
\item
  implementing a file naming strategy
\item
  backing up your work
\end{itemize}

Done consistently, and done well, RDM prevents data loss, helps ensure our work has long-term value and allows for the possibility of collaboration and sharing.

\chapter{Data Management Planning}\label{data-management-planning}

Why plan? Beyond all the benefits a good plan provides a research team, some funders require it.

\begin{itemize}
\tightlist
\item
  Tri-Agency Research Data Management Policy, Released March 15, 2021.
\item
  Tri-Agency Statement of Principles on Digital Data Management, Released June 15, 2016.
\item
  Tri-Agency Open Access Policy on Publications, Last modified December 21, 2016.
\item
  Access to Research Results: Guiding Principles
\end{itemize}

\chapter{Data Security}\label{data-security}

\end{document}
